%%%PAGE 121 rot=0 legibility=good summary=天体物理：万有引力定律发现史与推导(开普勒→牛顿→月地检验)、非球体万有引力计算(割补法，两幅图与公式)
%%%BLOCK id=121-01 ch=05 sec=万有引力定律·发现史 kind=knowledge alt=none cont=none y=0.05-0.245
% 页面右上角露出邻页(化学)内容“ρ=1.429g·L^-1”，不属于本页，已忽略
\textbf{天体物理}

托勒密\ 地心说\quad／\quad 哥白尼\ 日心说\quad／\quad 第谷\ 观测\quad／\quad 开普勒在第谷基础上总结开三\quad／\quad 伽利略发现木卫

胡克、笛卡尔、牛顿 $\to$ 万有引力定律

$F_{\text{万}}=G\dfrac{m_1m_2}{r^2}$\quad（\unclear{天体}球心距）

$G=6.67\times10^{-11}\ \unit{N\cdot m^2/kg^2}$\quad 卡文迪许扭称实验测出（放大法） % CHECK: 原稿“扭称”，应为“扭秤”
%%%END
%%%BLOCK id=121-02 ch=05 sec=万有引力定律·推导与月地检验 kind=derivation alt=04 cont=none y=0.245-0.43
第谷 $\to$ 开普勒三大定律 $\to$ 胡克：平方反比.（引力与$r^2$成反比）$\to$ 牛顿\quad 由牛一.\ 行星必受力

由牛二\quad $F_{\text{万}}=\dfrac{mv^2}{r}$

\[
\begin{aligned}
&\text{由圆周}\ v=\frac{2\pi r}{T}\\
&\text{由开三}\ \frac{r^3}{T^2}=\text{定值}
\end{aligned}
\qquad
\left.\begin{aligned}
F_{\text{万}}&=m\frac{4\pi^2}{T^2}r\\
T^2&=\frac{r^3}{\text{定值}}
\end{aligned}\right\}
\ \to\ F_{\text{万}}=\frac{m\,4\pi^2}{r^2}\cdot\text{定值}
\]

\[ F_{\text{万}}\propto\frac{1}{r^2},\ \propto m\ \ \propto M \]

$\downarrow$

\[ F_{\text{万}}=\frac{GMm}{r^2}\qquad \text{\unclear{对}月地检验}\ \to\ \text{正确}\qquad \frac{g}{g}=\frac{1}{60^2} \]% CHECK: “月地检验”前一字难辨(疑为“对”)，原稿此处有带箭头的下划线指向“正确”；g/g 的分子分母疑带撇或下标，原稿难辨
%%%END
%%%BLOCK id=121-03 ch=05 sec=非球体万有引力计算 kind=method alt=none cont=none y=0.40-0.96
\textbf{非球体万有引力计算}

割补法、对称填补法、负质量法

对剩余部分对$m$万有引力

%FIG p121_a 0.045 0.495 0.625 0.668
\notefig[0.58]{figures/p121_a.png}
% 图中标注：大球 M、半径 R，挖去部分(红色斜线小球)“挖 1/8 M”；红字“补”；F_1(红色箭头)、F_2；d-R/2、d、d-1/2R 等尺寸线；中间小圆圈为质点 m
% 图外：右侧小球右上方另有红字“补”
\red{补}

%FIG p121_b 0.04 0.665 0.33 0.935
\notefig[0.3]{figures/p121_b.png}
% 图中标注：余弦/正弦定理、(合力)F_2、质点 m、d、R、斜边 $\sqrt{d^2+(R/2)^2}$；红色箭头；红色斜线为挖去的小球

\begin{align*}
F_1&=\frac{G\frac{1}{8}Mm}{\left(d-\frac{1}{2}R\right)^2}\\
F_1+F_2&=\frac{GMm}{d^2}\\
F_{\text{万}}&=F_r+F_2-F_1 \\
&=\frac{GMm}{d^2}-\frac{G\frac{1}{8}Mm}{\left(d-\frac{1}{2}R\right)^2}
\end{align*}% CHECK: F_万=F_r+F_2-F_1 中第一项下标原稿像 r，疑为 1
%%%END
%%%PAGEEND 121
